\chapter{Reproductibilité, déploiement et exploitation}

\section{Principe de reproductibilité}

La reproductibilité repose sur quatre éléments liés : une entrée identifiée, un protocole paramétré, un \texttt{run\_id} et un artefact de sortie. Une valeur du mémoire n’est considérée comme finale que si cette chaîne peut être reconstruite. Les tableaux agrégés facilitent la lecture, mais les JSON et CSV bruts restent la preuve principale.

Les manifests conservent les chemins, tailles et empreintes SHA-256 lorsque celles-ci sont disponibles. Cette empreinte prouve que deux exécutions utilisent le même fichier ; elle ne suffit pas à établir qu’une copie locale est une distribution officielle si la provenance n’est pas documentée.

\section{Organisation des preuves}

Les expériences finales sont réparties par phase :

\begin{itemize}
  \item \texttt{experiments/phase\_dataset\_strengthening} contient les protocoles HDFS, BGL, multi-format et CSE-CIC-IDS2018 ;
  \item \texttt{experiments/phase\_multi\_agent} conserve le benchmark A--D et la campagne multi-VM ;
  \item \texttt{experiments/phase\_final\_scientific\_consolidation} rassemble les contrôles forensiques, le bootstrap, le routeur open-set et l’exécution multi-source avec modèles réels ;
  \item \texttt{revision\_reports} contient l’audit du manuscrit, le catalogue et le générateur des figures finales.
\end{itemize}

Dans chaque phase, \texttt{raw} conserve la sortie la plus proche de l’exécution, \texttt{processed} les représentations intermédiaires, \texttt{aggregated} les synthèses et \texttt{figures} les visualisations. Les graphiques insérés dans le mémoire sont copiés dans \texttt{figures/final} pour stabiliser la compilation.

\section{Chaîne de reproduction}

La reproduction suit le même ordre pour toutes les campagnes. Il faut d’abord vérifier l’entrée et son hash, installer les dépendances figées, exécuter le script avec les paramètres et graines enregistrés, puis comparer le schéma et les valeurs de l’artefact. La figure est régénérée seulement après cette comparaison.

Le script \texttt{revision\_reports/generate\_final\_scientific\_figures.py} lit les artefacts validés et applique un style commun. Aucune métrique expérimentale n’est saisie manuellement lorsqu’elle existe dans un CSV ou un JSON. Les titres, unités, tailles de police et sources sont centralisés afin que les figures restent lisibles après réduction A4.

\section{Paramètres déterminants}

Pour CICIDS2017, la reproduction doit préciser la séparation : aléatoire, scénario ou temporelle. Les graines 42 à 46 sont utilisées dans les comparaisons principales. La table des cinq candidats repose sur un espace commun et un holdout fichier/scénario.

Pour HDFS, les paramètres essentiels sont les blocs ordonnés par première apparition, les proportions 60/20/20, les échantillons 6,000/2,000/2,000, la graine 42 et Drain3 avec similarité 0,4, profondeur 4 et 100 enfants. L’agrégateur moyenne et le seuil 1,2729429339548877 sont gelés après validation.

Pour BGL, l’artefact final doit conserver le seuil Histogram, le groupe de connaissance du template et le baseline de nouveauté. Pour CSE-CIC-IDS2018, il faut préserver les journées 15 et 16 février, les 78 caractéristiques, les cinq graines et les échantillons équilibrés de 20,000 lignes.

Le routeur utilise 31 fichiers indépendants, 28 connus et 3 inconnus. Le seuil 100 est sélectionné sur les connus seulement. Changer ce seuil après lecture des inconnus produirait une autre expérience.

\section{Reproduction du benchmark multi-agent}

Le benchmark A--D doit conserver quatre charges, quatre architectures et dix répétitions. Les 160 exécutions représentent 264,000 tâches. Les métriques sont calculées avec leurs unités natives : tâches/s, millisecondes, secondes CPU, pourcentage équivalent cœur, pourcentage normalisé machine et mégaoctets RSS.

L’environnement influence fortement les micro-latences. Une reproduction sur une autre machine doit donc comparer les tendances relatives et publier sa propre configuration, plutôt que d’exiger une égalité bit à bit des durées. Le résultat fonctionnel attendu reste l’absence d’échecs si les mêmes chemins sont couverts.

\section{Reproduction de la campagne multi-VM}

La campagne de référence est \texttt{multivm\_cnp\_20260909T202632Z\_37592}. Elle nécessite un coordinateur, Redis 7.4.9 et deux agents enregistrés : \texttt{debian-alpha} et \texttt{ubuntu-beta}. Les profils autorisent un traitement parallèle par agent. Le scénario contient 60 tâches, un échec contrôlé, une réattribution et un replay idempotent.

La preuve attendue inclut les hôtes et versions de système, les compteurs de messages, les résultats par agent et la section \texttt{idempotence}. Une reproduction valide ne dépend pas d’un partage strict 30/30, car l’ordonnancement peut varier ; elle doit en revanche terminer les tâches, observer la reprise injectée et ne pas dupliquer l’effet logique.

Les identifiants de connexion et mots de passe ne sont pas intégrés au mémoire ni aux artefacts publiables. Ils doivent être fournis par des variables d’environnement ou un gestionnaire de secrets. Le Redis de laboratoire reste centralisé et isolé ; l’ouverture réseau non authentifiée n’est pas un scénario de déploiement recommandé.

\section{Compilation du mémoire}

Le document utilise pdfLaTeX avec encodage UTF-8, police T1 et support français. La bibliographie est construite entre les passes LaTeX, puis deux passes supplémentaires stabilisent les références, la table des matières et les listes. Toutes les compilations finales sont lancées depuis le dossier \texttt{ARIEL\_LOGMINER\_MEMOIRE\_FINAL}.

Les contrôles portent sur les références non résolues, les débordements significatifs, les labels dupliqués, les caractères corrompus et la présence des figures. Le PDF n’est retenu qu’après stabilisation du nombre de pages.

\section{Exploitation locale}

Un déploiement local peut utiliser \texttt{LocalMessageBus} et \texttt{SQLiteIdempotencyStore}. Cette configuration convient aux tests et à l’analyse sur une seule machine. Le passage à Redis permet plusieurs processus et machines, mais ajoute des exigences de disponibilité, d’authentification et de supervision.

L’analyste consulte les résultats, vérifie la provenance, examine les événements corrélés et formule un retour. Une anomalie candidate n’est pas automatiquement un incident confirmé. Le dashboard facilite cette lecture, tandis que les artefacts restent la référence pour un audit détaillé.

\section{Frontière d’industrialisation}

Une mise en production demanderait des travaux non évalués ici : TLS, gestion de secrets, contrôle d’accès, journalisation inviolable, réplication Redis, tests de partition, sauvegarde, reprise après sinistre, surveillance continue et politique de rétention. Une orchestration de conteneurs pourrait faciliter le déploiement, mais n’améliorerait pas à elle seule la validité des modèles.

La montée en charge devrait être testée avec des inférences réalistes, sur une durée longue et avec des pannes du transport. Les objectifs de latence et de disponibilité devraient être fixés avant la campagne. Tant que ces éléments manquent, Ariel Logminer reste un prototype de laboratoire.

\section{Checklist de vérification}

Avant publication, les points suivants sont contrôlés :

\begin{itemize}
  \item fichiers texte encodés en UTF-8 et absence de caractères corrompus ;
  \item correspondance des valeurs entre texte, tableaux, figures et artefacts ;
  \item séparation explicite des datasets, splits et unités ;
  \item labels LaTeX uniques et références résolues ;
  \item figures lisibles sur A4 avec axes, unités et source ;
  \item bibliographie compilée et accents conservés ;
  \item limites industrielles maintenues dans les perspectives ;
  \item absence de donnée d’accès ou de secret dans la livraison.
\end{itemize}
